% =====================================================================
\part{Literature Review}
% =====================================================================

\section{Introduction}\label{sec:intro}

\paragraph{Three notions of causality.}
When a computer scientist says that one part of a system \emph{causes} another, they may mean at least three distinct things. In Pearl's theory of causal inference \cite{pearl2009causality}, $X$ has a causal effect on $Y$ if \emph{intervening} on $X$---forcing it to a value, regardless of how it would otherwise have been produced---changes the distribution of $Y$. In quantum foundations and in the theory of higher-order quantum processes \cite{kissinger2019categorical,oreshkov2012noncausal}, the causal structure of a process is given by its \emph{signalling} constraints, which inputs can influence which outputs, together with the further question of whether it is compatible with a definite causal order between its parts. In probabilistic modelling, a causal structure is a directed acyclic graph along which a joint distribution \emph{factorises}, and the content of the structure is the family of conditional independences it entails \cite{pearl2009causality,fritz2023dseparation}.

Three further usages should be named so as to be set aside. In categorical quantum mechanics, a process is called \emph{causal} when it preserves the discarding map \cite{PQP_Coecke_Kissinger_2017}. This is the causality axiom of Chiribella, D'Ariano and Perinotti \cite{chiribella2010purification}, glossed as ``no signalling from the future'', and Coecke \cite{coecke2014terminality} shows that it implies the no-signalling principle for processes composed along a causal structure; it is the base axiom on which signalling structure is built, not a notion of causality between parts. Fritz \cite{fritz2020synthetic} calls a Markov category \emph{causal} when the output of a morphism cannot depend on a choice made downstream of it; this is again a property of the ambient category. In artificial intelligence and in verification, \emph{actual causality} \cite{halpern2016actual} asks which particular event caused which particular outcome in a single run; that question is not treated here. Nor are experimental tests of indefinite causal order, causal discovery, or identifiability, each of which is mentioned below only where a work package touches it.

That the same word carries these meanings is not an accident, but neither is it a theorem. Each notion is \emph{extensional}: it characterises causal structure by what an observer may do to the system---intervene, send signals, condition. What such observers probe is \emph{intensional}: the dataflow of a program, the temporal order of its communication events, the graph of a model. Programs carry this intensional structure explicitly in their syntax. The viewpoint of this thesis is therefore that causality is a property of \emph{program semantics}, and that the right way to relate the three notions is to relate, for each one, the intensional structure in the syntax to the extensional characterisation in the semantics by a correspondence theorem.

\paragraph{Where the two routes stand.}
Two lines of work precede this proposal, and the review is organised around them. On the interventional side, my coauthors and I built a hierarchy of first-order probabilistic languages in which Pearl's interventions and counterfactual \emph{twinning} are program transformations, defined program equivalences contextually, and gave a graded denotational semantics that is fully abstract for concrete models (\S\ref{sec:intv}). This route has the complete picture for one notion of causality: the intensional dataflow of a program is characterised extensionally by equivalence under all interventional contexts. On the communication side, my MSc dissertation gave an algebraic theory $\iotq$ of classically controlled quantum communication, an operational model in which programs exchange qubits with a quantum stream, and a denotational model given by a free I/O monad on $\CPinf$, with the denotational model adequate and fully abstract for the operational one; this year's work, with my supervisor, is its conference write-up and a draft proof of completeness (\S\ref{sec:order}). Here the temporal order of communication events is first-class semantic data: the denotation is indexed by I/O traces. This route does \emph{not} yet have an extensional causal reading of that order; giving it one, as signalling structure in the sense of higher-order quantum processes, is part of the research plan rather than a result.

A single example shows why these are two faces of one question. In the interventional language, a program reads an \emph{intervention label} $\ell$ through $\mathsf{applyIntv}\ \ell\ t$, which yields the label's value when it is set and $t$ otherwise; if a second variable $z$ does not depend on the result, every intervention at $\ell$ leaves the distribution of $z$ unchanged, and the contextual equivalence makes this an equation between programs. Causal independence is \emph{tested by contexts}. In the theory $\iotq$, the commutation axiom $\applyU{U}(\vec a,\ \vec a.\ \opin(b.\ x(\vec a, b))) = \opin(b.\ \applyU{U}(\vec a,\ \vec a.\ x(\vec a, b)))$ says that a local operation and a communication event on different qubits may be reordered, and the quantum-stream semantics validates it; by contrast $\opin$ and $\opout$ do not commute with one another, and that is exactly why the trace index is a total order. Causal independence is \emph{an equation validated by an observer}. Both statements say that no information flows from one place to another, so that reordering or intervening is unobservable; in programming-language terms both are \emph{noninterference} properties \cite{goguen1982noninterference}. Whether this coincidence is an instance of a general correspondence is the question of Part II.

\paragraph{Scope of the review.}
Sections \ref{sec:foundations}--\ref{sec:order-wherewestand} treat the literature on which the thesis builds: the semantics of effects, probabilistic programs and quantum programs; causal models, their equivalences and their diagrammatic presentations; causal probabilistic programming; and quantum process calculi. Sections \ref{sec:concurrency} and \ref{sec:hoqc} position two neighbouring fields that the research plan reaches into, the semantics of concurrency and higher-order quantum processes, at the level of their main definitions and results; their detailed survey is a milestone of that plan (Part II, \S\ref{sec:plan}). Section \ref{sec:synthesis} draws the three notions together and identifies the gap the thesis addresses.

% ---------------------------------------------------------------------
\section{Semantic foundations shared by both routes}\label{sec:foundations}

\subsection{Monads, algebraic theories, and their combinations}

Moggi's observation that notions of computation are monads \cite{MOGGI199155,moggi1991monadicmetalanguage} and Plotkin and Power's refinement that most of them are presented by algebraic theories of operations and equations \cite{plotkin_power_2002,PlotkinPower04CompEffectsOverview} are the common foundation of both routes. An algebraic theory determines a monad, and conversely a finitary monad is a Lawvere theory \cite{lawveretheories_monads_HYLAND2007437}; the equational presentation is what makes program equivalence a matter of derivability. Hyland, Plotkin and Power \cite{sums_tensors_HYLAND200670} showed that the standard ways of combining theories have computational meaning: the \emph{sum} adjoins the operations of two theories with no interaction, so that the order of effects from the two theories is observable, while the \emph{tensor} additionally imposes that operations of one theory commute with those of the other, so that their relative order is unobservable. In the premonoidal semantics of effects \cite{premonoidal_cat_POWER_ROBINSON_1997,PowerRobinson97premonoidal,fuhrmann1999premonoidal}, the morphisms that commute with everything are the \emph{central} ones; centrality is the established algebraic form of ``cannot be observed to happen before or after''. Handlers \cite{bauer2019algebraicalgebraiceffectshandlers} give the programming counterpart of models; comodels and runners \cite{model_comodel_PLOTKIN2008295,runners_10.1007/978-3-030-44914-8_2} describe the environments in which effectful programs execute.

Ordinary algebraic theories cannot express operations that create and consume \emph{resources}, such as qubits, which by no-cloning must be treated linearly. Staton's \emph{parameterised algebraic theories} (PATs) \cite{statonlics13,statonpopl15} extend the framework with parameters: operations consume parameters, classically choose a branch, and bind fresh parameters in each branch. Staton used them to give a fully complete equational theory of qubit quantum computation \cite{statonpopl15}; Stark's free model of $\pi$-calculus name creation \cite{stark:freamp} is the earlier instance of the same commutation discipline, where communication primitives commute with independent local operations. Kammar, Liell-Cock, Lindley, Matache and Staton \cite{kammar2026threads} have given a complete equational axiomatisation of dynamic threads (fork, wait, stop, atomic actions) as a parameterised algebraic theory, whose closed terms correspond exactly to pomsets; I return to this in \S\ref{sec:concurrency}, since it is the nearest neighbour of my own completeness result.

\emph{Where we stand.} My MSc dissertation \cite{wang_algebraic_2024} shows that sums and tensors exist for the enriched Lawvere theories presented by PATs, by explicit syntactic construction, which is what licenses the definition of $\iotq$ as a tensor in \S\ref{sec:order}.

\emph{What this gives us.} The tensor, and centrality more generally, are already primitive statements about causal independence. In $\iotq$ the tensor makes communication commute with local operations on other qubits, while $\opin$ and $\opout$ do not commute with each other. I treat this as an observation to be made precise, not as a result: Part II asks (\WP{3}) whether centrality in an algebraic theory is in general the algebraic shadow of no-signalling.

\subsection{Semantics of probabilistic programs}

Probabilistic programming gives a compositional syntax for statistical models and performs Bayesian inference on them compositionally \cite{vandemeent2021introductionprobabilisticprogramming}. Its denotational semantics is monadic: the Giry monad on measurable spaces \cite{kock72,semantics_ppl_staton2016}, and, for higher-order programs, the probability monad on quasi-Borel spaces \cite{Heunen_2017}. Implementations such as LazyPPL \cite{lazyppl}, MonadBayes \cite{monadbayes1,monadbayes2} and GenJAX \cite{becker26genjax} realise this monadic structure directly. Full abstraction is the recurring theorem shape for such languages, the standard reference point being probabilistic coherence spaces for probabilistic PCF \cite{ehrhard2014pcoh}.

The synthetic approach of \emph{Markov categories} \cite{fritz2020synthetic,fritz2019probability} axiomatises the categorical structure needed to talk about probability without measure theory: a symmetric monoidal category with a commutative comonoid on every object, where discarding is natural (the unit is terminal) but copying need not be. In Markov categories with conditionals, disintegration and Bayesian inversion are expressed synthetically \cite{cho2019disintegration,fritz2020synthetic}, and Fritz and Klingler \cite{fritz2023dseparation} prove the soundness of d-separation, and the equivalence of the local and global Markov properties, synthetically; completeness of d-separation in the sense of Geiger, Verma and Pearl remains the classical result. This is the \emph{factorisation} face of causality: a Bayesian network is a string diagram in a Markov category, and the conditional independences it entails are read off the diagram.

The second piece of structure both routes rely on is \emph{grading}. Graded monads \cite{katsumata2014_graded_strong_monad,orchard2020unifying,fujii2016towards,dorsch2019graded} and, more generally, locally graded categories \cite{levy2019locallygraded,sarkis2025string} track an index---an effect annotation, a trace, a resource---through composition. Milius, Pattinson and Schröder \cite{milius2015gradedtrace} use graded monads for trace semantics; Liell-Cock and Staton's $\mathbf{ImP}$ construction \cite{liell2025compositional} uses grading to make imprecise probability compositional. The interventional semantics of \S\ref{sec:intv-wherewestand} is built on an analogous construction in which the grade records which interventions are available.

\emph{What this gives us and what it lacks.} Markov categories give factorisation but no primitive notion of intervention; the diagrammatic causal-model literature (\S\ref{sec:diagrammatic}) adds it by surgery on diagrams. Neither offers a notion of program \emph{context}, which is what contextual equivalence needs.

\subsection{Semantics of quantum programs}

Selinger's QPL \cite{QPL_SELINGER_2004} established the denotational template for quantum programming languages: types denote finite-dimensional spaces of density matrices (or, more generally, $C^*$-algebras \cite{SemOperatorAlgebra}), and programs denote completely positive maps. Categorical quantum mechanics \cite{abramsky2007categoricalsemanticsquantumprotocols,heunen_vicary_book,PQP_Coecke_Kissinger_2017} gave this a graphical calculus through dagger compact closed categories and the CPM construction \cite{selinger_dagger_compact_closed_cat_cp_maps}, and the ZX-family of calculi axiomatise fragments of it; Staton's PAT $\mathsf{Q}$ \cite{statonpopl15} and Carette et al.'s extension of $\Pi$ \cite{Quantum_PI_Carette_2024} give complete equational theories of qubit programs and of unitary circuits respectively. Languages such as QWIRE, Quipper and Silq \cite{qwire_10.1145/3009837.3009894,quipper_DBLP:journals/corr/abs-1304-3390,silq} work in the circuit model with classical control. For languages with infinitary features one needs the completion of $\CP$ under infinite biproducts, $\CPinf$ \cite{PaganiSV13,D_Completion_ZHAO20102167}, which is the base category of the communication semantics. Quantum monads \cite{sati2023quantummonadology,QIO_Altenkirch_Green_2009} and quantum game semantics \cite{clairambault2019quantumgames} are two further lines of work on the structure of quantum computation as an effect; the latter uses event structures and is the closest existing semantics to the one proposed in \WP{2a}.

\emph{What this gives us.} The axiom that singles out the deterministic (discard-preserving) processes in categorical quantum mechanics \cite{PQP_Coecke_Kissinger_2017,chiribella2010purification} is the base condition from which the signalling structure of \S\ref{sec:hoqc} is built; the $\CausC$ construction requires compact closure and further ``precausality'' axioms on top of it, and, by design, its base category also contains the non-deterministic, trace-decreasing processes.

% ---------------------------------------------------------------------
\section{Causality as intervention}\label{sec:intv}

\subsection{Structural causal models and Pearl's hierarchy}

Pearl's structural causal models (SCMs) \cite{pearl2009causality,Pearl2018_book_of_why} are systems of equations $X_i = f_i(\mathrm{Pa}_i, U_i)$ in which each endogenous variable is a deterministic function of its parents and of an exogenous noise variable; a prior on the noise makes an SCM a probabilistic model, and the equations make it more than a distribution. The surplus is organised by the \emph{causal hierarchy} \cite{bareinboim:etal20}: rung 1 asks observational questions; rung 2 asks \emph{interventional} questions, answered by replacing the equation for $X$ with a constant (the $\dok$ operator); rung 3 asks \emph{counterfactual} questions about a particular run, answered by abduction, action and prediction, implemented graphically by \emph{twinning}, which duplicates the model into worlds that share the noise \cite{Balke1994}. Both operations are \emph{syntactic transformations of the model}. The Neyman--Rubin potential-outcomes framework \cite{rubin1974} is the other dominant formalism; its correspondence with SCMs is well understood \cite{ibeling2023comparing}. Richardson and Robins' single-world intervention graphs \cite{richardson2013swig} present an intervention by splitting a node into its natural and its fixed value, which is the graphical counterpart of the intervention labels of our languages; they are deliberately single-world, and cross-world quantities such as natural direct effects need further assumptions there. Interventions need not be perfect: the literature on stochastic and soft interventions \cite{dawid2002_stoch_intv,didelez2006_stoch_mech_intv,Eberhardt_Scheines_2007,levis2024stochasticinterventionssensitivityanalysis} replaces an equation by another mechanism rather than a constant.

Throughout, I work with fully specified models and Bayesian inference: mechanisms and priors are given, and queries are answered by inference. \emph{Identifiability} (when an interventional quantity is computable from observational data, the subject of the do-calculus \cite{JMLR:v9:shpitser08a,shpitser2012counterfactualstested}) and \emph{causal discovery} \cite{spirtes2000causation,peters_janzing_schoelkopf_2017_book} are outside the scope of this thesis.

\emph{What this gives us.} The causal structure of an SCM is its dataflow: which variables each equation reads. Interventions probe it. The semantic question is when two models have the same causal structure, and the answer must be extensional.

\subsection{Equivalences of causal models and logical accounts}

Bongers, Forré, Peters and Mooij \cite{Bongers_2021} give the principal account of equivalence for SCMs at all three rungs: two models are observationally, interventionally or counterfactually equivalent when they agree on the corresponding class of queries, with the counterfactual case defined through the two-world twinned model. Their treatment also covers cyclic SCMs, through solution functions. At rung 2, several further notions exist: Pearl's causal Bayesian networks are families of interventional distributions compatible with a graph; \emph{causal spaces} \cite{park2024measuretheoreticaxiomatisationcausality,buchholz2024products} axiomatise this measure-theoretically; and \emph{causal abstraction} \cite{rubenstein2017consistency,beckers2019abstracting,rischel2021abstraction,lorenztull2026abstraction,englberger2025causalabstractionscategoricallyunified} studies coarsening maps between models that commute with interventions. At the level of graphs rather than models, Markov equivalence \cite{verma1990equivalence} is what discovery algorithms can recover. On the logical side, Ibeling and Icard \cite{ibeling2020hierarchy} characterise the three rungs by three probabilistic logics of strictly increasing expressivity whose semantics is defined on SCMs through Pearl's operations, and Icard has proposed probabilistic Turing machines as the computational counterpart of SCMs \cite{icard2017programs,ibeling2019openuniverse}.

\emph{What this gives us and what it lacks.} These are all extensional characterisations of interventional structure, and they are the yardstick against which our contextual equivalences are measured (\S\ref{sec:intv-wherewestand}). None of them is formulated as a contextual equivalence of programs, and none treats higher-order models.

\subsection{Diagrammatic and categorical causal models}\label{sec:diagrammatic}

Fong's \emph{causal theories} \cite{fong2012causaltheories} present a Bayesian network as a symmetric monoidal category freely generated by its graph, with models as functors into stochastic maps. Jacobs, Kissinger and Zanasi \cite{jacobs2019causal} present interventions as \emph{string-diagram surgery} in this setting, and introduce \emph{comb disintegration}: the use of combs, processes with a hole, to express the classical causal-inference computations, which is the first appearance of the comb notion of \S\ref{sec:hoqc} in classical causal inference. Lorenz and Tull \cite{lorenz2023causalmodelsstringdiagrams} develop causal models in string diagrams systematically, and a later line presents causal abstraction as natural transformations between such models \cite{lorenztull2026abstraction,englberger2025causalabstractionscategoricallyunified}. Bayesian inference in this style is handled by \emph{normalisation boxes}, which renormalise locally inside a diagram \cite{jacobsszelesstein2025inference}, and Agrawal et al.\ \cite{agrawal2026equational} obtain equational principles for causal Bayesian networks from them. Zhang et al.\ \cite{zhang2026semanticseparation} characterise d-separation as a noninterference-style separation property, which makes explicit a connection implicit in all of this work and in \S\ref{sec:intro}: invariance under interventions is a noninterference property \cite{goguen1982noninterference}.

\emph{What this gives us and what it lacks.} This is the closest existing bridge between the factorisation face and the intervention face of causality, and it is compositional in the functorial sense. Its equational theories are, however, stated per graph, and it has no notion of program context or of higher-order models. It is where the unifying work package (\WP{3}) starts, and comb disintegration is the prior work that \WP{3a} must be positioned against.

\subsection{Causal probabilistic programming languages}

Causal probabilistic programming aims to make intervention and twinning primitives of an expressive probabilistic language. Libraries such as DoWhy \cite{pywhy,blobaum2024dowhygcm} and the $\dok$ operators of Pyro and PyMC \cite{bingham2018pyro,abrilpla2023pymc} manipulate models as graphs or computation graphs and stop at rung 2 (DoWhy's GCM module reaches rung 3 on fitted mechanisms). Witty et al.\ \cite{witty2019synthesis,witty2023bayesian} treat an intervention as a source transformation that replaces subterms, which is powerful but non-compositional; Brulé's Whittemore \cite{brule2018whittemore} is an early embedded causal DSL.

Four systems reach rung 3 natively. OmegaC \cite{tavares2021language} obtains interventions by a lazy evaluation strategy that retroactively modifies closures, reminiscent of dynamic binding, and counterfactuals by making every program a deterministic function of explicit noise. ChiRho \cite{chirho} implements both as effect handlers over Pyro, with a multi-world value type whose relation to Pearl's twinning is not immediate. MultiVerse \cite{perov2020multiverse} uses trace inspection and re-execution. Kiesel et al.\ \cite{kiesel2023whatif} answer counterfactual queries on ProbLog programs by a transformation proved correct against Pearl's semantics.

\emph{What this gives us and what it lacks.} These systems show that the primitives are implementable at higher order, by non-standard evaluation techniques. What they lack is a denotational and equational foundation: a semantics in which intervention is a mathematical operation rather than an evaluation strategy, and an account of when two causal programs are equivalent. That is what the first route supplies.

\subsection{Where we stand on the interventional route}\label{sec:intv-wherewestand}

Our POPL submission \fillin{status: under revision for POPL 2027?; coauthors} (an early version appeared at the LAFI workshop \cite{lafi-wang}; artifact \cite{ARTIFACT_wang_2026_23207221}) reconstructs Pearl's hierarchy as a hierarchy of three first-order probabilistic languages. $\Lk{1}$ is a standard first-order probabilistic language with conditioning by a final score. $\Lk{2}$ exposes hooks in the dataflow as typed free variables, the \emph{intervention labels}, read through $\mathsf{applyIntv}$; an intervention $\dok(\ell\mapsto v)\{t\}$ is substitution of $v$ for $\ell$. $\Lk{3}$ annotates each binding as exogenous or endogenous, and \emph{twinning} $\mathsf{twin}_n$ is a vectorisation transformation from $\Lk{3}$ to $\Lk{2}$ that duplicates the endogenous bindings into $n$ worlds sharing the exogenous ones. Program equivalence at rung $k$ is \emph{contextual} with respect to rung-1 denotational equality, in Milner's sense: two programs are $\equiv_k$ when they are denotationally equal in every rung-$k$ context, and varying the admissible interventions and world counts gives a spectrum of equivalences. For representable acyclic SCMs compiled into $\Lk{3}$, the observational, interventional and counterfactual equivalences of Bongers et al.\ coincide with the equivalences $\equiv_3^{\mathrm{obs}}$, $\equiv_3^{\mathrm{intv}}$ and $\equiv_3^{\text{2-world}}$ of the compiled programs; the full $\equiv_3$ quantifies over all $n$-world twinnings, and whether it is strictly finer than $\equiv_3^{\text{2-world}}$ is open. The denotational semantics interprets interventions by \emph{grading}, in a construction reminiscent of $\mathbf{ImP}$ \cite{liell2025compositional}, and twinning by world-polymorphic vectorisation. It is adequate for the contextual equivalences unconditionally and fully abstract when the underlying model is \emph{concrete}, meaning that values correspond to global elements and global elements separate morphisms; $\mathbf{Set}$, standard Borel spaces and quasi-Borel spaces are concrete, and the construction is parametric in the model. The denotation then motivates a higher-order language $\LHO$ with cross-world interventions, into which $\Lk{3}$ compiles in a semantics-preserving way, and which has a definitional Haskell implementation. \fillin{one sentence on which parts are yours: e.g.\ the contextual-equivalence framework and the full-abstraction proofs}

The limitations relevant to Part II are these: the languages are acyclic; the relation between the contextual equational theory and the normalisation-box equational theories of \S\ref{sec:diagrammatic} is not worked out; the first-order semantics is over normalised probability, while $\LHO$ widens to the unnormalised measure monad to internalise conditioning, so that programs with observations denote sub-normalised maps; and the models are measure-theoretic rather than finite-dimensional, which matters for the signalling face (\S\ref{sec:synthesis}).

% ---------------------------------------------------------------------
\section{Causality as order and signalling in communicating quantum programs}\label{sec:order}

\subsection{Quantum communication and distributed quantum computing}

Quantum communication underlies teleportation \cite{Teleportation_PhysRevLett.70.1895}, key distribution \cite{BB84_BENNETT20147} and, increasingly, proposals for distributed quantum computing that connect small quantum processors into larger ones \cite{DQC_survey_Caleffi_2024}. Programming models for distributed quantum computing are at present circuit-level: they describe which gates are applied where and which qubits are sent when. On the semantic side, Ceragioli, Gadducci, Lomurno and Tedeschi observe that the existing quantum process calculi do not agree on what the observational properties of quantum values should be \cite{lqCCS_10.1145/3632885}. There is, in other words, no agreed account of what a distributed quantum program \emph{means}, nor of which of its behaviours are physically realisable.

\subsection{Quantum process calculi}

Quantum communication has mostly been studied through quantum process calculi---CQP \cite{CQP_gay2004communicatingquantumprocesses}, QPAlg \cite{QPAlg_lalire2004processalgebraicapproachconcurrent}, qCCS \cite{qCCS_ying2010algebraquantumprocesses} and lqCCS \cite{lqCCS_10.1145/3632885}---which extend classical calculi \cite{ccs_10.5555/539036,CSP_10.1145/359576.359585,pi_calc_MILNER19921} with qubit message passing. They are presented operationally, by labelled transition systems over configurations, with program equivalence given by bisimilarity. As the authors of lqCCS show, the proposed bisimilarities differ in subtle ways because of the interaction of concurrency, synchronisation and local nondeterminism with quantum effects; their later coalgebraic account of effect-labelled transition systems \cite{ceragiolicoalgebraic} is an attempt to organise these subtleties. A separate difficulty is well known from the classical theory: parallel composition is not an algebraic operation and cannot be adjoined to an algebraic theory as an operation \cite{stark:freamp,CSP_van_Glabbeek_2010}, although workarounds exist, such as the cooperative threads of Abadi and Plotkin \cite{Cooperative_Threads_DBLP:journals/corr/abs-1009-2405}.

\emph{What this gives us and what it lacks.} Process calculi are the only existing semantic treatment of communicating quantum programs, and they are expressive. They are operational only, their equivalences are not canonical, and they entangle communication with concurrency, which makes it hard to isolate what is specific to \emph{communication}. My work studies communication in isolation, as an algebraic effect, and is in that sense the local sublanguage of these calculi with a denotational semantics.

\subsection{Where we stand on the communication route}\label{sec:order-wherewestand}

My MSc dissertation \cite{wang_algebraic_2024}, and its conference version with my supervisor \fillin{FoSSaCS submission status}, introduce $\iotq$, an algebraic theory of classically controlled qubit communication, as the \emph{tensor} of a standard I/O theory \cite{sums_tensors_HYLAND200670} with Staton's complete theory $\mathsf{Q}$ \cite{statonpopl15}: programs may send and receive qubits depending on measurement outcomes, and communication commutes with local quantum operations. The tensor is argued to be needed, since one equation that holds in the tensor but not in the sum is required to validate a quantum phenomenon involving entanglement across the program--environment boundary (Example B in Part II). The theory is given two models. The operational model lets a program exchange qubits with an environment modelled as a \emph{quantum stream}; a program equivalence $\approx_q$ compatible with quantum theory is defined from it by inspecting final quantum states, and terms modulo $\approx_q$ form a model of $\iotq$, which shows that the tensor does not collapse in the manner of Eckmann--Hilton. Although operational, $\approx_q$ is therefore not a contextual equivalence; whether a contextual equivalence observing only traces, continuations and a final measurement coincides with it is a question I return to in Part II. The denotational model is the free I/O monad $T(X)=\bigoplus_{\sigma}\mathrm{qbit}^{|\sigma|}\otimes X$ on $\CPinf$, indexed by I/O traces $\sigma$. The main theorem is that the denotational semantics is adequate and fully abstract for the operational one: $\llbracket c\rrbracket = \llbracket c'\rrbracket$ iff $c\approx_q c'$. What the dissertation left as a conjecture is \emph{completeness}: that $\llbracket t\rrbracket=\llbracket t'\rrbracket$ implies $t=t'$ is derivable in $\iotq$. I now have a draft proof of this (manuscript in preparation, not yet closed), by induction on the I/O depth of terms, using a reorganisation step that normalises every term into an I/O-free head applied to a single input and a single output branch, and a factorisation lemma on completely positive maps that lets semantically equal terms be rewritten to a common prefix. If it closes, the three semantics---algebraic, operational, denotational---coincide.

The observation that sets up Part II is this. The trace $\sigma$ is a \emph{temporal} order of communication events, exposed in the denotation: the $\sigma$-component of a closed term is a completely positive map with $|\sigma|$ open qubit wires, one per event, to be used in the order $\sigma$, and when the term branches on measurements this component is sub-normalised, its trace being the probability of the branches that produce $\sigma$. Whether the family of these components is a \emph{causal} object in the signalling sense of \S\ref{sec:hoqc}, a comb or a tester whose slots are ordered by $\sigma$, is a question the dissertation raised and did not answer. It is stated as Conjecture \ref{conj:b1} in Part II.

\subsection{Order in concurrency theory}\label{sec:concurrency}

Concurrency theory has long modelled causal order explicitly. Lamport's happens-before relation \cite{lamport1978time} is the partial order an observer of a distributed system can reconstruct. Mazurkiewicz traces \cite{diekert1995traces} are words modulo the commutation of independent letters, the cheapest partial-order model and the one that arises when independent channels are added to a sequential program. Pratt's pomsets \cite{pratt1986pomsets}---partially ordered multisets of events, up to isomorphism, with series and parallel composition---generalise words to partially ordered words, and Gischer \cite{gischer1988pomsets} axiomatised the equational theory of series-parallel pomsets. Winskel's event structures \cite{winskel1986eventstructures,nielsen1981petri,winskel1995models} add a conflict relation to the causal order and give a domain-theoretic account of nondeterministic concurrent computation. Concurrent Kleene algebra \cite{hoare2011cka} adds a commutative parallel operator to Kleene algebra, related to sequential composition by an exchange law; Kappé, Brunet, Silva and Zanasi \cite{kappe2018cka} prove that series-parallel pomset languages closed under Gischer's subsumption order are its free model, by reduction to an earlier completeness theorem for bi-Kleene algebra \cite{laurence2014completeness}. Within the algebraic-effects tradition, Dvir, Kammar and Lahav \cite{dvir2022sharedstate} give a monadic semantics of shared-state concurrency from an equational theory, and Kammar et al.\ \cite{kammar2026threads} axiomatise dynamic threads as a parameterised algebraic theory whose closed terms are exactly pomsets, with a completeness theorem. Clairambault, de Visme and Winskel \cite{clairambault2019quantumgames} give an event-structure game semantics for quantum programs.

\emph{Position.} The trace-indexed monad of \S\ref{sec:order-wherewestand} is the degenerate single-channel case of this picture: a single channel forces a total order, so traces are words. Two extensions are distinguished in Part II. Adding independent channels to a \emph{sequential} program makes events on different channels commute, so that traces become Mazurkiewicz traces and the monad a free monad over a partially commutative monoid; this keeps the proof architecture of the single-channel case. Adding \emph{threads} makes traces pomsets, and the theory of Kammar et al.\ is a parameterised algebraic theory with exactly that free model, so ``tensor $\iotq$ with a thread theory'' is a well-posed, and much harder, second step. Completeness results in this literature also go through free models; whether their techniques transfer to the quantum setting, where a factorisation lemma on CP maps did the work, is part of the first reading milestone of the plan. The logics of session types are the other point of contact: BV and its relatives have been used as logics of multiparty protocols, and Simmons conjectures that causal consistency of a wiring of black boxes corresponds to deadlock freedom of the protocol \cite{simmons2024thesis}; a trace-indexed semantics of communicating programs is where such a correspondence could be made precise.

\subsection{Higher-order quantum processes and signalling}\label{sec:hoqc}

Higher-order quantum computation studies transformations of quantum processes rather than of states. Chiribella, D'Ariano and Perinotti's \emph{quantum combs} \cite{quantum_circuit_architecture_Chiribella_2008,chiribella2009networks} are processes with holes in a fixed causal order; their probabilistic counterparts, \emph{testers}, and Gutoski and Watrous's \emph{strategies} \cite{gutoski2007games} describe multi-round interactions with sub-normalised outcomes, and the general theory of supermaps forms a hierarchy of types \cite{bisio2019hoqt}. Oreshkov, Costa and Brukner's \emph{process matrices} \cite{oreshkov2012noncausal} are the most general bipartite process operators compatible with quantum mechanics holding locally in each party's laboratory; the multipartite theory and the notion of \emph{causal separability}, in which the order of later parties may depend classically on earlier outcomes, are due to Araújo et al.\ and Oreshkov and Giarmatzi \cite{araujo2015witnessing,oreshkov2016separable}. Some process matrices are causally non-separable. Two kinds of evidence distinguish them: device-independent violation of a \emph{causal inequality}, which Oreshkov, Costa and Brukner's process exhibits, and device-dependent witnesses of non-separability. The \emph{quantum switch} \cite{chiribella2013switch}, which applies two operations in an order controlled by a qubit in superposition, is causally non-separable \cite{araujo2015witnessing} but produces only causal correlations; photonic experiments implement it \cite{procopio2015experimental}, and whether such implementations realise or simulate indefinite order is debated \cite{oreshkov2019timedelocalized,wechs2023timedelocalised}.

Which indefinite-order processes can be built as circuits is the subject of Wechs, Dourdent, Abbott and Branciard \cite{wechs2021qcqc}, who define nested classes of quantum circuits with fixed order, with \emph{classical} control of causal order (the next operation may depend on classical outcomes of earlier ones) and with \emph{quantum} control of causal order, and characterise each by conditions on process matrices. The switch lies in the quantum-control class and not in the classical-control one; every classical-control process is causally separable; and quantum-control processes generate only causal correlations. That the quantum-control class exhausts the physically realisable processes is a proposal, not a theorem.

Categorically, Kissinger and Uijlen's $\CausC$ construction \cite{kissinger2019categorical} builds, over a \emph{precausal} category---compact closed, with a chosen discarding structure and further axioms, the examples being $\CP$ on finite-dimensional systems and $\Mat$, finite matrices of non-negative reals---a $*$-autonomous category whose types are signalling constraints: first-order processes, non-signalling processes, one-way signalling processes, and so on. The dual of the non-signalling type contains all bipartite process matrices, causally separable or not; separability is a property inside that type, not a type, and Simmons proves that every process matrix is an affine combination of definite-order processes \cite{simmons2024thesis}. Simmons's thesis \cite{simmons2024thesis} is the fullest development of this line and the reference for most of what Part II needs on the signalling side. It rebuilds the construction over \emph{additive} precausal categories (compact closed with products, hence biproducts, and a finite causal basis for every object), which gives $\CausC$ products, coproducts and probabilistic choice as type operators; it adds a non-commutative sequential operator for one-way signalling, under which $\CausC$ models Guglielmi's BV-logic \cite{simmons2022bv}; it defines $n$-ary \emph{graph types} encoding compatibility with an arbitrary DAG of signalling relations; it surveys candidate ways of admitting sub-normalised, partial processes into the construction, effectus-style, and where each falls short; and it introduces \emph{causal logic}, pomset logic with directed axiom links, whose proof nets characterise exactly the causally consistent wirings of black boxes \cite{simmons2024thesis,simmons2024completelogic}, with BV and pomset logic separated by a process-matrix example. The finite-causal-basis axiom excludes infinite-dimensional and measure-theoretic bases, which is one of the two obstructions discussed in \S\ref{sec:synthesis}. Finally, \emph{quantum causal models} transfer Pearl's framework to quantum processes: Costa and Shrapnel \cite{costa2016quantumcausalmodelling} gave the first process-matrix-based definition, Allen et al.\ \cite{allen2017quantumcommoncauses} the common-cause line, and Barrett, Lorenz and Oreshkov \cite{barrett2019quantumcausalmodels} the general framework, in which a node is an input--output pair, the process operator factorises as a product of commuting operators (a quantum Markov condition), interventions are local instruments at nodes, and d-separation in the graph implies conditional independence of classical outcomes. The cyclic extension \cite{barrett2021cyclic} gives a cyclic-causal-model account of some processes without definite order, including the switch. Coherent control of channels \cite{coherent_control_channel_Abbott_2020,resource_theory_comm_Kristjánsson_2020} and, on the language side, the PBS-calculus of Clément and Perdrix \cite{clement2020pbs}, which has the switch as a term, are the relevant prior work on coherent control.

\emph{Position.} Three things follow for this thesis. First, the family of $\sigma$-components of an $\iotq$ denotation has the shape of a tester or strategy whose slots are ordered by the traces; that it \emph{is} one is Conjecture \ref{conj:b1}. Second, signalling types are the extensional face of order in the same way that interventional equivalence is the extensional face of dataflow, and quantum causal models are the one existing place where factorisation, intervention and signalling are all present, which is why they are the first reading milestone. Third, the analogy between parallel composition and indefinite causal order must be stated with care, and the process-matrix picture is the right way to state it: the parties' local programs are combs, and the process matrix is the \emph{network} that wires them together, in other words the scheduler. Parallel composition $P\parallel Q$ leaves the order of events unspecified. A classical scheduler, even one whose later choices depend on earlier outcomes, yields a process in the classical-control class, hence a causally separable one. The switch is a coherent scheduler, well-typed on every non-signalling pair of channels, whose output applies the two operations in a superposition of orders; it is causally non-separable, and it has content only when the two programs act on a shared target or communicate, since on disjoint systems every order gives the same channel. The analogy is therefore between $\parallel$ and the argument type of a scheduler, not between $\parallel$ and the switch itself; this is what makes \WP{2d} well-posed, and also what makes it hard, since coherent control of order is not an algebraic operation with classical continuations. Finally, Simmons's own outlook \cite{simmons2024thesis} proposes reading send and receive operations as physical processes so as to assign causal types to communicating programs, with causal consistency playing the role of deadlock freedom and the session-type logics BV, GV and MAV as the point of contact; that is, from the physics side, the direction of \WP{2a} and \WP{2b}.

% ---------------------------------------------------------------------
\section{Synthesis: three notions of causality and the gap}\label{sec:synthesis}

Table \ref{tab:synthesis} summarises the review. Each column is one extensional notion of causality; each row asks the same question of it.

\begin{table}[ht]
\centering\small
\begin{tabularx}{\textwidth}{>{\raggedright\arraybackslash}p{3.1cm}XXX}
\toprule
 & \textbf{Intervention} & \textbf{Signalling} & \textbf{Factorisation}\\
\midrule
Intensional datum & dataflow of a program / SCM & temporal order of events: trace $\sigma$, comb order, pomset & DAG\\
Extensional characterisation & contextual equivalence under interventions & signalling types ($\CausC$); causal (non-)separability & conditional independence / d-separation\\
Correspondence theorem & full abstraction for concrete models (\S\ref{sec:intv-wherewestand}) & full abstraction and draft completeness relate three equalities (\S\ref{sec:order-wherewestand}); the signalling reading is Conjecture \ref{conj:b1} & soundness of d-separation, synthetically \cite{fritz2023dseparation}; completeness classically\\
Semantics of programs & first-order and higher-order, acyclic & single channel, classical control & Markov categories; interventions by surgery\\
Next hard case & cyclic SCMs & multiple channels; causally non-separable processes & cyclic and quantum causal models\\
\bottomrule
\end{tabularx}
\caption{Three notions of causality and where the two routes stand.}\label{tab:synthesis}
\end{table}

Three remarks. First, the example of \S\ref{sec:intro} and the asymmetry it exposes: for intervention the correspondence between intensional and extensional structure is a theorem; for signalling there is a semantics with order as first-class data and a conjecture. Second, non-DAG structure is the next hard case in every column, but it is not the same case: feedback in cyclic SCMs, causal non-separability in quantum processes, and unscheduled parallel composition in programs are three different departures from a fixed order, and a unifying account must distinguish them. Third, and this is the gap: there is no framework in which the three columns are structure on one category. Markov categories cover factorisation and, through surgery, part of intervention; $\CausC$ covers signalling for quantum and finite classical processes; quantum causal models relate all three, with a compositional presentation in Lorenz and Tull's string diagrams, but are not formulated as a semantics of programs. The two semantics of \S\ref{sec:intv-wherewestand} and \S\ref{sec:order-wherewestand} cover the first two columns separately, and both are families of morphisms indexed by an intensional datum, intervention grades or traces, with a reindexing action, which is the common structure a generalisation would start from.

Two obstructions should be displayed rather than hidden. The first is \emph{finite dimension}: the $\CausC$ construction needs a compact closed base with invertible dimensions, which $\Mat$ and finite-dimensional $\CP$ provide, but measurable spaces, quasi-Borel spaces and the infinite-biproduct completion $\CPinf$ do not. The bridge of \WP{3a} is therefore stated on the finite fragment, and the comb reading of \WP{2b} is stated per trace, over finite-dimensional $\CP$, with the trace index kept external. The second is \emph{sub-normalisation}: conditioning in $\LHO$ denotes sub-normalised maps, and branching in $\iotq$ makes individual trace components sub-normalised, so neither lands in a first-order $\CausC$ type as it stands; testers and strategies \cite{chiribella2009networks,gutoski2007games} are the objects that accommodate sub-normalised families, Simmons \cite{simmons2024thesis} surveys the ways of admitting partial processes into $\CausC$ and finds each wanting in some respect, and normalisation boxes \cite{jacobsszelesstein2025inference} are the diagrammatic tool for renormalising after conditioning. I also note that the factorisation column is the thinnest for me: it enters through Markov categories and the diagrammatic causal-model literature, not through my own results, and \WP{3a} includes a factorisation reading for that reason.

